\documentclass[10pt,a4paper]{amsart}
\usepackage[margin=19mm]{geometry}
\usepackage{amsmath,amssymb,mathtools}
\usepackage[scr=rsfs,cal=euler]{mathalfa}
\usepackage{xfrac}
\usepackage[unicode,hidelinks,bookmarks=false]{hyperref}
\newcommand{\ie}{i.e.\ }
\newcommand{\cf}{cf.\ }
\newcommand{\resp}{respectively\ }
\newcommand{\Subsection}{\subsection}
\providecommand{\nobreakdash}{\nobreak\hbox{-}}
\setlength{\emergencystretch}{2em}
\multlinegap0pt
\usepackage{bm}
\usepackage{bbm}
\usepackage{dsfont}
\usepackage{xspace}
\usepackage{cancel}
\usepackage{mathtools}
\usepackage{amsthm}
\makeatletter
\@ifundefined{theorem}{\newtheorem{theorem}{Theorem}}{}
\makeatother
\title[Nonlinear Fractal Histopolation Function]{Nonlinear Fractal Histopolation Function}
\author{Aiswarya T, Srijanani Anurag Prasad}
\date{Source: arXiv:2509.18657v1 (CC BY 4.0)}
\begin{document}
\maketitle

\noindent\emph{Source and licence.} Nonlinear Fractal Histopolation Function. Version arXiv:2509.18657v1 (CC BY 4.0), distributed under the Creative Commons Attribution 4.0 International licence, as recorded in the arXiv licence metadata for this exact version and shown on the abstract page https://arxiv.org/abs/2509.18657v1. Full rights notice: licenses/arxiv-2509.18657-CC-BY-4.0.txt.\par\smallskip

\noindent\textit{Editorial scope.} Counted unit: the Theorem whose source label is \texttt{graphoffp} in the article (source0.tex lines 184--196, source label \texttt{graphoffp}), delivered together with the article's own complete original proof of it (source0.tex lines 197--267, one \texttt{proof} environment of 71 lines). No mathematics is written, repaired or re-derived: statement and proof are the article's own text line for line. Required context retained uncounted: l\_j (lines 123--125); F\_j (lines 129--131); RB (lines 142--144). Every definition, display and label of those lines is retained unchanged. Omitted from this package and not redistributed: 1--122, 126--128, 132--141, 145--183, 268--590. Nothing inside a retained excerpt is omitted or altered: every retained body is delivered exactly as the article's lines are written, so no omission receipt of any kind accompanies this package. Citations: the retained excerpts carry no citation command at all, so no bibliography entry of the article is transcribed and no reference is redistributed. Reference adaptation: none was needed and none was made; every reference inside the delivered unit resolves inside it, so no unresolved reference or citation is produced. Printed slips kept literal: no printed word, symbol, hypothesis or number of the counted statement or of its proof was altered. Layout and class: the article's own class is \texttt{sn-jnl} with packages including graphicx, multirow, amsmath, amssymb, amsfonts, amsthm, mathrsfs, appendix, xcolor, textcomp, manyfoot, booktabs, algorithm, algorithmicx; the wrapper is the lane's pinned \texttt{amsart} editorial wrapper at 10pt on A4, which restates the theorem-like environments the retained text needs and adds the article's own macro definitions that the retained text needs (none), with extra wrapper packages (bm, bbm, dsfont, xspace, cancel, mathtools). The delivered numbering is the unit's own. Assets: no retained span contains a figure environment, a graphics-inclusion command, a PDF-page inclusion, a table or a TikZ picture; no image file is copied into this package, the package declares no asset dependency and no image file is redistributed with it. Licence: version arXiv:2509.18657v1 of this article is distributed under the Creative Commons Attribution 4.0 International licence, as recorded in the arXiv licence metadata for this exact version and shown on the abstract page https://arxiv.org/abs/2509.18657v1. A full rights notice is delivered at licenses/arxiv-2509.18657-CC-BY-4.0.txt. Characters: every retained excerpt is pure ASCII, so the delivered engine, XeLaTeX with its default Latin Modern text font, reports no missing character.
\begin{equation}
l_j(t)=a_jt+b_j  \label{l_j}
\end{equation}

\begin{equation}
F_j(t,x)=c_jt+\delta_js_j(x)+d_j, \label{F_j}
\end{equation}

\begin{eqnarray}
Th(t)&=F_j(l_j^{-1}(t),h(l_j^{-1}(t))) = c_jl_j^{-1}(t)+\delta_js_j(h(l_j^{-1}(t)))+d_j\label{RB}.
\end{eqnarray}

\begin{theorem}\label{graphoffp}
Let $F_j$ and $l_j$ be as defined in Eq.~\eqref{F_j} and Eq.~\eqref{l_j}. Consider~
\\$\mathcal{I}=\left\{ A; w_j, j\in\mathbb N_N\right\}$ where  $w_j: A\rightarrow A $ are defined by 
\begin{equation}
w_j(t,x)=(l_j(t),F_j(t,x)).             \label{wi}
\end{equation}
Let $C=\max\limits_{j\in\mathbb N_N}|c_j|$, where $c_j$ are constants given in Eq.~\eqref{F_j}. Then, each of the maps $w_j$ is a Rakotch contraction with respect to the metric $d_{\eta}$ given as
\begin{equation*}
d_{\eta}((t,x),(t',x')):=|t-t'|+\eta|x-x'|,
\end{equation*}
where $\eta=\frac{1-\max\limits_{j\in \mathbb N_N}|a_j|}{2(C+1)}$.
Further, a unique attractor exists for $\mathcal{I}$. The attractor is the closure of the graph of the fixed point of the R.B. operator given in Eq.~\eqref{RB}.
\end{theorem}

\begin{proof}
Consider $(t,x),(t',x') \in I\times\mathbb{R}$. Then
\begin{eqnarray*}
    |F_j(t,x)-F_j(t',x')|&=&|c_jt+\delta_js_j(x)+d_j-(c_jt'+\delta_js_j(x')+d_j)|\\
    &\leq&C|t-t'|+\delta|s_j(x)-s_j(x')|\leq C|t-t'|+\zeta(|x-x'|),
\end{eqnarray*}
where $\zeta(t)=\delta~\psi(t)$.
\begin{eqnarray*}
d_\eta(w_j(t,x),w_j(t',x'))&=& d_\eta((l_j(t),F_j(t,x)),(l_j(t'),F_j(t',x')))\\
&=& |l_j(t)-l_j(t')|+\eta|F_j(t,x)-F_j(t',x')|\\
&\leq&(|a_j|+\eta C)|t-t'|+\eta\zeta(|x-x'|)\\
&\leq&\bigg(|a_j|+\eta C+\eta\frac{\zeta(|t-t'|+|x-x'|)}{|t-t'|+|x-x'|}\bigg)|t-t'|\\
& &\hspace{2cm}+\eta\frac{\zeta(|t-t'|+|x-x'|)}{|t-t'|+|x-x'|}(|x-x'|).\\
%&\leq&(|a_j|+\eta C+\eta)|x_-x'|\\
%& &\hspace{2cm}+\eta\frac{\psi(|x-x'|+|y-y_1|)}{|x-x'|+|y-y_1|}(|y-y_1|)\\
\end{eqnarray*}
As $\frac{\zeta(s)}{s}<1 $ for all $ 
s>0,$
\begin{eqnarray*}
d_\eta(w_j(t,x),w_j(t',x'))&\leq&(\max\limits_{j\in\mathbb N_N} |a_j|+\eta C+\eta)|t-t'| \\
& &\hspace{2cm}+\eta\frac{\zeta(|t-t'|+|x-x'|)}{|t-t'|+|x-x'|}(|x-x'|)\\
&\leq&\max\{\max\limits_{j\in\mathbb N_N
} |a_j|+\eta C+\eta,\frac{\zeta(|t-t'|+|x-x'|)}{|t-t'|+|x-x'|}\} \times \\
& &\hspace{5cm} d_\eta((t,x),(t',x')).
\end{eqnarray*}
 Since $\eta=\frac{1-\max\limits_{j\in\mathbb N_N}|a_j|}{2(C+1)}$, $\max\limits_{j\in\mathbb N_N}|a_j|+\eta C+\eta<1$. Note that $\frac{\zeta(s)}{s}<1$ for $s>0$. Thus,
\begin{equation*}
\max\{\max\limits_{j\in\mathbb N_N} |a_j|+\eta C+\eta,\frac{\zeta(|t-t'|+|x-x'|)}{|t-t'|+|x-x'|}\}<1.
\end{equation*}
For $s>0$, define
\begin{equation*}
\gamma(s):=\max\{\max\limits_{j\in\mathbb N_N} |a_j|+\eta C+\eta,\frac{\zeta(s)}{s}\}.
\end{equation*}
Since $\gamma(s)$ is a non-increasing function,
\begin{eqnarray*}
d_\eta(w_j(t,x),w_j(t',x'))&\leq& \gamma(d((t,x),(t',x')))d_\eta((t,x),(t',x'))\\
&\leq&\gamma(d_\eta((t,x),(t',x')))d_\eta((t,x),(t',x')).
\end{eqnarray*}
Thus, $w_j$ are Rakotch contractions, and $\mathcal{I}$ has a unique attractor, say $B$.

Let $\mathcal{H}(A)$ be the space comprising nonempty compact subsets of $A$ equipped with the Hausdorff metric.
Define $W:\mathcal{H}(A)\rightarrow\mathcal{H}(A)$ as
\begin{equation*}
W(K)=\bigcup\limits_{j=1}^N w_j(K).
\end{equation*}
Then, $B=W(B)$. 

Let $f$ be the fixed point of the R.B. operator given by Eq.~\eqref{RB} and $G$ be its graph, i.e.,
\begin{equation*}
G=\left\{(t,f(t)):t\in I\right\}.
\end{equation*}
Let $(t,x)\in \bar{G}$. So, there exists a sequence $t_m\in I$ such that $(t_m,f(t_m))\in G$ converges to $(t,x)$. Since $I$ is a disjoint union of $I_j$, there exists a unique $j_m\in\mathbb N_N$ such that $t_m\in I_{j_m}$. So $t_m=l_{j_m}(\tilde{t_m})$ for some $\tilde{t_m}\in I$. Further,
\begin{eqnarray*}
(t_m,f(t_m))=(l_{j_m}(\tilde{t_m}),f(l_{j_m}(\tilde{t_m})))
=(l_{j_m}(\tilde{t_m}),F_{j_m}(\tilde{t_m},f(\tilde{t_m})))=w_{j_m}(\tilde{t_m},f(\tilde{t_m})),
\end{eqnarray*}
 which implies $(t_m,f(t_m))\in W(G)\subset W(\bar{G})$. Thus, $(t,x)\in W(\bar{G})=\overline{W(
G )}$. So, 
\begin{equation}
\bar{G}\subseteq W(\bar{G}).   \label{eq:G1}   
\end{equation}
Conversely, let $(t,x) \in W(\bar{G})=\overline{W(G)}$. Now, there exists a sequence $\{(t_m,x_m)\} \in W(G)$ such that $(t_m,x_m) \rightarrow (t,x)$ as $m \rightarrow \infty$. Since $(t_m,x_m) \in W(G)$, there exists a unique $j_m\in\mathbb N_N$ such that 
\begin{eqnarray*}
(t_m, x_m) = w_{j_m}(\tilde{t_m}, f(\tilde{t_m})) = (l_{j_m}(\tilde{t_m}),F_{j_m}(\tilde{t_m},f(\tilde{t_m}))) = (l_{j_m}(\tilde{t_m}),f(l_{j_m}(\tilde{t_m}))).  
\end{eqnarray*} 
This implies that $(t_m,x_m) = (l_{j_m}(\tilde{t_m}),f(l_{j_m}(\tilde{t_m}))) \in G$ and hence $(t,x) \in \bar{G}$. This gives
\begin{equation}
W(\bar{G})\subseteq \bar{G}. \label{eq:G2}   
\end{equation}
Eq.~\eqref{eq:G1} and Eq.~\eqref{eq:G2} together implies that $W(\bar{G})=\bar{G}$. Since the attractor $B$ is unique which satisfy the equation  $B=W(B)$, $B=\bar{G}$. 
\end{proof}	

\end{document}
